\documentclass[sigplan,screen,review,anonymous]{acmart}

\setcopyright{none}
\acmConference[POPL 2027]{ACM SIGPLAN Symposium on Principles of Programming Languages}{January 2027}{TBD}
\acmYear{2027}
\acmDOI{}
\acmISBN{}

\usepackage{booktabs}
\usepackage{listings}
\usepackage{xspace}

\newcommand{\Agda}{Agda\xspace}
\newcommand{\Qupit}{\textsc{qupit}\xspace}
\newcommand{\Zp}{\mathbb{Z}_p}
\newcommand{\Word}{\mathsf{Word}}
\newcommand{\Gen}{\mathsf{Gen}}
\newcommand{\NF}{\mathsf{NF}}
\newcommand{\LM}{\mathsf{LM}}
\newcommand{\ABox}{\mathsf{A}}
\newcommand{\BBox}{\mathsf{B}}
\newcommand{\DBox}{\mathsf{D}}
\newcommand{\EBox}{\mathsf{E}}
\newcommand{\LBox}{\mathsf{L}}
\newcommand{\MBox}{\mathsf{M}}
\newcommand{\file}[1]{\texttt{#1}}
\newcommand{\Sp}{\mathrm{Sp}}
\newcommand{\SL}{\mathrm{SL}}
\newcommand{\GL}{\mathrm{GL}}
\newcommand{\Cliff}{\mathcal{C}}
\newcommand{\PauliGrp}{\mathcal{P}}
\newcommand{\Zps}{\Zp^{\ast}}

\lstdefinelanguage{AgdaLite}{
  morekeywords={module,where,open,import,data,record,let,in,with,private,variable,using,renaming,hiding,public},
  sensitive=true,
  morecomment=[l]{--}
}
\lstset{
  language=AgdaLite,
  basicstyle=\ttfamily\small,
  columns=fullflexible,
  keepspaces=true,
  frame=single,
  breaklines=true
}

\begin{document}

\title{A Mechanized Account of the \Qupit Agda Development}
\subtitle{Deriving Box Relations for Odd-Prime-Dimensional Clifford Circuits}

\author{Anonymous Author(s)}
\affiliation{%
  \institution{Anonymous Institution}
  \country{}
}

\begin{abstract}
This paper explains the \Qupit \Agda repository that accompanies the
paper ``A Complete and Natural Rule Set for Multi-Qudit Clifford
Circuits in All Odd Prime Dimensions.''  The development formalizes
Clifford circuit syntax as words over generators, equips these words
with presentation-generated congruence relations, defines modular
arithmetic over finite fields \(\Zp\), and uses those ingredients to
derive box-style rewrite relations for one-, two-, and three-qupit
fragments.  The repository is not merely a transcription of circuit
diagrams: its core abstractions separate raw word syntax, presentation
closure, algebraic normal forms, Pauli action, and checked equational
derivations.  This separation is what makes the proof scripts usable:
large circuit equalities are decomposed into finite-field arithmetic,
associativity normalization, lifting lemmas for extra wires, and
small reusable Clifford identities.

The paper is written as an artifact-oriented account of the code.  It
identifies the trusted and untrusted boundaries in the development,
summarizes the implemented relation sets, explains how boxes are
encoded and interpreted back into circuits, and maps the public
box-relation statements in \file{BoxRelations.agda} to their proof
modules under \file{N/BR}.  Several older or exploratory modules in
the repository contain unfinished holes; following the artifact itself,
we treat those modules as outside the described proof spine.
\end{abstract}

\ccsdesc[500]{Theory of computation~Program verification}
\ccsdesc[500]{Theory of computation~Equational logic and rewriting}
\ccsdesc[300]{Hardware~Quantum computation}

\keywords{Agda, Clifford circuits, qudits, presentations, mechanized mathematics, rewriting}

\maketitle

\section{Introduction}

Reasoning about Clifford circuits is algebraic, but the algebra is
large enough that informal derivations become brittle.  A typical box
relation states that a primitive gate can be pushed through a
parameterized circuit pattern, producing another parameterized pattern
and a residual gate sequence.  Even when the high-level rule is
conceptually simple, a proof may require case analysis over finite
field elements, repeated use of modular arithmetic, and long
associativity rearrangements of circuit words.  The \Qupit repository
uses \Agda to make these derivations checkable.

The development targets Clifford circuits over odd prime dimensions.
The prime is represented in the code by parameters
\texttt{p-2 : Nat} and \texttt{p-prime : Prime (2+ p-2)}; the actual
dimension is \(p = 2 + \texttt{p-2}\).  The finite field carrier
\(\Zp\) is implemented with \texttt{Fin p}, and all circuit
parameters used by the box definitions are elements of this finite
type.  This choice lets the development use ordinary dependent types
for field-indexed circuits while keeping the prime assumption explicit
in every module that depends on it.

At the top level, the README says that the repository accompanies a
paper on a complete and natural rule set for multi-qudit Clifford
circuits.  The existing documentation, \file{doc/Readme.pdf}, gives a
compact explanation of the central proof idea: instead of doing the
box-relation calculations by hand, define the relevant presentations
inside \Agda and let the type checker validate the derivations.  This
paper expands that documentation into a POPL-style artifact paper,
with an emphasis on the engineering and proof architecture visible in
the code.

The most important constraint on this account is fidelity to the
repository.  The code base contains more than one layer of work:
generic presentation infrastructure, active Clifford and box-relation
modules, simplified and derived relation sets, and older experiments.
Some experimental modules contain unfinished holes.  We therefore
describe the implemented proof spine and explicitly separate it from
unused or unfinished material.

\paragraph{Contributions.}
This paper makes four contributions.

\begin{itemize}
\item It gives a module-by-module account of the \Qupit
  mechanization, from words and congruence closures to Clifford
  generators and box relations.
\item It explains why the repository represents circuits as binary
  words rather than lists, and how associativity is recovered by a
  presentation-level congruence and list-based normalization lemmas.
\item It documents the relation between the comprehensive
  \file{N/Symplectic.agda} relation set, the simplified
  \file{N/Symplectic-Simplified.agda} relation set, and the derived
  relation infrastructure used by Pauli actions.
\item It maps the public theorem index in \file{BoxRelations.agda} to
  the one-, two-, and three-qupit proof modules under \file{N/BR}.
\end{itemize}

\section{Mathematical Background}
\label{sec:math}

This section records the mathematics that the development mechanizes.
None of it is new; it fixes notation and explains why a presentation of
words, quotiented by a congruence, is an adequate model of Clifford
equality in odd prime dimension.  Throughout, \(p\) is an odd prime and
\(\Zp = \mathbb{Z}/p\mathbb{Z}\) is the field with \(p\) elements.

\subsection{The qudit Pauli group}

Fix a primitive \(p\)-th root of unity \(\omega = e^{2\pi i/p}\) and let
the shift and clock operators act on \(\mathbb{C}^p\) by
\[
  X\,|j\rangle = |j+1 \bmod p\rangle,
  \qquad
  Z\,|j\rangle = \omega^{\,j}\,|j\rangle .
\]
They satisfy \(X^p = Z^p = I\) and the Weyl commutation relation
\(ZX = \omega\,XZ\).  For \((a,b)\in\Zp^2\) the Weyl operator
\(W(a,b) = \omega^{-2^{-1}ab}\,X^a Z^b\) is well defined because \(2\) is
invertible in \(\Zp\); this is one of the places where oddness of \(p\)
is essential~\cite{Appleby2005,Gross2006}.  The \(n\)-qudit Pauli (or
Heisenberg--Weyl) group \(\PauliGrp_n\) is generated by the operators
\(W(a,b)\) on each wire together with the central phases
\(\langle\omega\rangle\).  Two Weyl operators commute up to a phase
governed by the symplectic form: writing \(v=(a,b)\) and
\(v'=(a',b')\),
\[
  W(v)\,W(v') = \omega^{\langle v,v'\rangle}\, W(v')\,W(v),
  \qquad
  \langle v,v'\rangle = a'b - a b' .
\]
This pairing is exactly \texttt{sform1} of \file{N/Pauli.agda} up to the
sign convention~\cite{Gottesman1999,Hostens2005}.

\subsection{The Clifford group and its symplectic quotient}

The \(n\)-qudit Clifford group \(\Cliff_n\) is the normalizer of
\(\PauliGrp_n\) in the unitary group \(U(p^n)\), taken modulo global
phase.  Conjugation by a Clifford operator sends Weyl operators to Weyl
operators, hence induces a \(\Zp\)-linear map on the label space
\(\Zp^{2n}\).  Because conjugation preserves Weyl commutators, that map
preserves \(\langle\,\cdot\,,\cdot\,\rangle\): it is symplectic.  This
is the content of the short exact sequence
\[
  1 \longrightarrow \PauliGrp_n/\langle\omega\rangle
    \longrightarrow \Cliff_n/\langle\omega\rangle
    \longrightarrow \Sp(2n,\Zp)
    \longrightarrow 1 .
\]
For odd \(p\) the sequence splits: every symplectic matrix is realized
by a Clifford unitary, the splitting being the Weil (metaplectic)
representation over a finite
field~\cite{Neuhauser2002,Appleby2005,Gross2006}.  Consequently, up to
Paulis and phases, the Clifford group \emph{is} \(\Sp(2n,\Zp)\).  This
is the structural fact that justifies the design of the repository: the
symbolic congruence on circuit words is sound for Clifford equality
precisely because both sides descend to the same symplectic matrix, and
completeness amounts to the surjectivity
above~\cite{Hostens2005,deBeaudrap2013,Farinholt2014}.

\subsection{The symplectic group over \(\Zp\)}

Let \(\Omega\) be the standard symplectic form on \(\Zp^{2n}\).  The
symplectic group is
\[
  \Sp(2n,\Zp)
    = \{\, M \in \GL_{2n}(\Zp) \mid M^\top \Omega\, M = \Omega \,\},
\]
and its order is
\[
  |\Sp(2n,\Zp)| = p^{\,n^2}\prod_{i=1}^{n}\bigl(p^{2i}-1\bigr).
\]
For \(n=1\) the symplectic condition reduces to \(\det M = 1\), so
\(\Sp(2,\Zp)=\SL_2(\Zp)\), of order \(p(p^2-1)\)~\cite{LidlNiederreiter}.
Finite presentations of these groups, of the kind the development
reconstructs at the level of circuit words, are classical in
combinatorial and computational group
theory~\cite{MagnusKarrassSolitar,Johnson1997,Sims1994}.

\subsection{Generators as symplectic matrices}

Write each transformation as a matrix acting on column vectors
\((a,b)^\top\), matching the one-qupit actions of \file{N/Action.agda}.
The phase gate, Hadamard, and multiplier act by
\[
  \rho(S^k) = \begin{pmatrix}1&0\\ k&1\end{pmatrix},
  \quad
  \rho(H) = \begin{pmatrix}0&-1\\ 1&0\end{pmatrix},
  \quad
  \rho(M(x)) = \begin{pmatrix}x^{-1}&0\\ 0&x\end{pmatrix},
\]
each of determinant \(1\), hence elements of \(\SL_2(\Zp)\).  The
one-qupit relations of \file{N/Symplectic.agda} are then literally
matrix identities over \(\Zp\):
\[
\begin{aligned}
  \rho(S)^p &= I, &\qquad
  \rho(H)^4 &= I, &\qquad
  (\rho(S)\rho(H))^3 &= I,\\
  \rho(H)^2 &= \rho(M(-1)) = -I, &
  \rho(M(x))\,\rho(M(y)) &= \rho(M(xy)), &
  \rho(M(x))\,\rho(S)\,\rho(M(x))^{-1} &= \rho(S)^{x^2}.
\end{aligned}
\]
For instance \(\rho(H)^2 = -I = \mathrm{diag}(-1,-1) = \rho(M(-1))\)
recovers the simplified axiom \(H^2 = M(-1)\), and the conjugation
identity recovers \(M(x)S = S^{x^2}M(x)\).  On two wires \(\rho(CZ)\) is
the symplectic shear
\[
  \rho(CZ):\ (a,b),(a',b') \longmapsto (a,\,b+a'),\,(a',\,b'+a),
\]
which preserves the form and commutes with sufficiently shifted
single-wire gates---exactly the structural relations the code records as
the Selinger relations~\cite{Selinger2015}.

\subsection{Boxes as symplectic coordinates}

The box families \(\ABox,\BBox,\DBox,\EBox\) and the normal forms
\(\NF(n)\) of \file{N/Boxes.agda} form a coordinate system for elements
of the symplectic group: each box stores the field parameters of one
elementary symplectic factor, and a normal form \(\NF(n)\) is an ordered
product of such factors, one layer per wire block.  A box relation
\([x]\cdot g \approx \mathit{dir}\cdot[x']\) then says that multiplying a
normal-form factor by a generator \(g\) can be re-expressed as a
direction word \(\mathit{dir}\) times an updated factor \([x']\); read
through \(\rho\), it asserts the equality of the two symplectic matrices.
Pushing a generator past the factors of a normal form and recomputing
the parameters by finite-field arithmetic is therefore a constructive
form of the Bruhat-style factorization of
\(\Sp(2n,\Zp)\)~\cite{Sims1994,HoltEickOBrien}.  This is the geometric
meaning of the otherwise syntactic derivations in \file{N/BR}.

\section{Repository Overview}

Table~\ref{tab:modules} summarizes the active parts of the project
that this paper discusses.  The table follows the structure of the
repository rather than imposing a separate mathematical organization.

\begin{table*}
\centering
\caption{Main modules used by the \Qupit development.}
\label{tab:modules}
\begin{tabular}{ll}
\toprule
File or directory & Role \\
\midrule
\file{Word/Base.agda} & Binary word syntax, powers, word maps, folds, and word relations. \\
\file{Word/Properties.agda} & Word-map laws and decidable equality for words. \\
\file{Presentation/Base.agda} & Congruence closure of a relation on words. \\
\file{Presentation/Properties.agda} & Setoid, magma, semigroup, monoid, and associativity tools. \\
\file{Presentation/Morphism*.agda} & Morphisms and identity/isomorphism infrastructure for presentations. \\
\file{Presentation/CosetNF.agda} & Generic coset and Reidemeister-Schreier-style normal-form machinery. \\
\file{Zp/ModularArithmetic.agda} & Modular arithmetic over \texttt{Fin n}; ring and field-like lemmas. \\
\file{Zp/Fermats-little-theorem.agda} & Prime-modulus and primitive-root support. \\
\file{N/Symplectic.agda} & Main generator set and comprehensive Clifford relations. \\
\file{N/Symplectic-Simplified.agda} & Paper-facing simplified relation set parameterized by a primitive root. \\
\file{N/Symplectic-Derived.agda} & Derived generators such as parameterized powers of \(H,S,CZ\). \\
\file{N/Iso-Sym-Derived.agda} & Isomorphism between the basic and derived symplectic presentations. \\
\file{N/Action.agda}, \file{N/Pauli.agda} & Pauli vectors, symplectic form, and generator action. \\
\file{N/Boxes.agda}, \file{N/LM-Sym.agda} & Box parameter spaces and interpretation of boxes as circuit words. \\
\file{N/BR/...} & Concrete box-relation derivations and supporting calculations. \\
\file{BoxRelations.agda} & Public index of box-relation statements and proof pointers. \\
\bottomrule
\end{tabular}
\end{table*}

The generic presentation layer is intentionally independent of
Clifford circuits.  It can be read as a small library for monoid
presentations: a presentation is a generator type \(X\), a relation
\(\Gamma : \Word X \to \Word X \to \mathsf{Set}\), and the congruence
closure of \(\Gamma\).  The Clifford-specific layer instantiates this
infrastructure with generators \(H\), \(S\), \(CZ\), and lifting
constructors that move generators to later wires.

\section{Words Rather Than Lists}

The basic syntax of circuits is defined in \file{Word/Base.agda}.
The definition is small but central:

\begin{lstlisting}
data Word (X : Set) : Set where
  [_]w : X -> Word X
  epsilon : Word X
  _bullet_ : Word X -> Word X -> Word X
\end{lstlisting}

In the actual code the constructor names use Unicode:
\texttt{[\_]\^{}w}, \texttt{epsilon} is written \(\epsilon\), and
\texttt{\_bullet\_} is written \(\_\bullet\_\).  The paper uses ASCII
names in listings where that improves portability.

The key design choice is that \(\Word X\) is a binary tree.  A list of
gates would quotient away the parse tree of composition too early:
the expression \(u \cdot v\) carries the fact that it was built from
left operand \(u\) and right operand \(v\), while the flattened list
of gates does not.  Many equational proofs need to rewrite only the
left or right side of a product; keeping the tree makes such proof
steps structurally visible to \Agda.

Associativity is not built into the syntax.  Instead, it is added as
a proof rule in the congruence closure.  This is a familiar tradeoff
in mechanized algebra: raw syntax stays simple and inductive, while
the semantic equality relation supplies associativity, units, and
domain axioms.

The word module also defines powers:
\[
  w^0 = \epsilon,\qquad
  w^1 = w,\qquad
  w^{n+2} = w \cdot w^{n+1}.
\]
Powers are used throughout the Clifford development for relations
such as \(S^p = \epsilon\), \(CZ^p = \epsilon\), and \(H^4 =
\epsilon\), and also for parameterized circuits such as \(S^k\) and
\(CZ^k\).

\section{Presentation Closure}

The congruence closure is implemented in \file{Presentation/Base.agda}.
For a relation \(\Gamma\) on words, the module defines an equality
\(\approx\) generated by reflexivity, symmetry, transitivity,
congruence, monoid laws, and the axioms in \(\Gamma\):

\begin{lstlisting}
data _approx_ : WRel X where
  refl       : w approx w
  sym        : w approx v -> v approx w
  trans      : w approx v -> v approx u -> w approx u
  cong       : w approx w' -> v approx v' ->
               w bullet v approx w' bullet v'
  assoc      : (w bullet v) bullet u approx w bullet (v bullet u)
  left-unit  : epsilon bullet w approx w
  right-unit : w bullet epsilon approx w
  axiom      : Gamma w v -> w approx v
\end{lstlisting}

The closure is the object that most later proofs inhabit.  A theorem
such as ``pushing \(S\) through an \(\EBox\) box'' is an \Agda term
whose type is an equality in this closure.  Setoid reasoning then
provides the readable calculation style.

The module \file{Presentation/Properties.agda} packages the closure
as an \Agda setoid and proves that word composition forms a magma, a
semigroup, and a monoid modulo \(\approx\).  It also supplies a
practical associativity layer.  The functions \texttt{to-list} and
\texttt{from-list} flatten and rebuild words, and the lemma
\texttt{by-assoc} proves that equal flattened lists imply equality of
the original words modulo associativity.  Later proof scripts use this
infrastructure, directly or through tactics in
\file{Presentation/Tactics.agda}, to avoid drowning in parenthesis
management.

\section{Finite Fields}

The finite-field layer lives under \file{Zp}.  The carrier of
\(\mathbb{Z}/n\mathbb{Z}\) is \texttt{Fin n}:

\begin{lstlisting}
Z : Nat -> Set
Z n = Fin n
\end{lstlisting}

Nonzero elements are represented as dependent pairs:

\begin{lstlisting}
Z* : Nat -> Set
Z* zero    = Z zero
Z* (suc n) = Sigma[ a in Z (suc n) ] (a /= zero)
\end{lstlisting}

Addition and multiplication are implemented by converting to natural
numbers, applying ordinary addition or multiplication, and reducing
modulo the modulus.  The file then proves the expected algebraic laws:
associativity and commutativity of addition and multiplication,
identities, additive inverse, distributivity, and additional lemmas
about powers and inverses.  For the Clifford development, the prime
modulus modules expose the dimension
\[
  p = 2 + \texttt{p-2}
\]
and field-specific operations such as inverse on nonzero elements.

This layer matters because the box relations are parameterized by
field elements.  For example, an \(\ABox\) contains a nonzero pair
\((a,b)\), an \(\EBox\) contains one field element, and the
interpretation of boxes uses expressions such as \(-b/a\).  In the
code, \(-b/a\) is not a partial operation; it is introduced only in
branches where the denominator is accompanied by a proof that it is
nonzero.

\section{Clifford Generators}

The main generator type is defined inside the module
\file{N/Symplectic.agda}.  For \(n\) qupits, generators are indexed by
the number of wires on which they are well-formed:

\begin{lstlisting}
data Gen : Nat -> Set where
  H-gen  : {n : Nat} -> Gen (1+n)
  S-gen  : {n : Nat} -> Gen (1+n)
  CZ-gen : {n : Nat} -> Gen (2+n)
  _lift  : {n : Nat} -> Gen n -> Gen (suc n)
\end{lstlisting}

The real code writes \texttt{\_lift} as a Unicode up-shift
constructor.  The lifting constructor shifts a generator onto later
wires.  Thus the first-wire gates are named directly, and gates on
later wires are obtained by repeated lifting.  Words are lifted with a
map:
\[
  w^\uparrow = \mathsf{wmap}\;(\_\!\uparrow)\;w.
\]
The code also defines a down-shift operation, but in the active
\file{N/Symplectic.agda} presentation \texttt{\_down} is implemented
as the identity on words.  Several older comments show a more
structural down-shift version, but the proof spine described here
uses the current code.

The primitive circuits are:
\[
S = [S\text{-gen}],\quad
H = [H\text{-gen}],\quad
CZ = [CZ\text{-gen}].
\]
From these, the repository defines common derived circuits:
\[
\begin{aligned}
S^{-1} &= S^{p-1},\\
CZ^{-1} &= CZ^{p-1},\\
CX &= H^\downarrow{}^3 \cdot CZ \cdot H^\downarrow,\\
XC &= H^\uparrow{}^3 \cdot CZ \cdot H^\uparrow.
\end{aligned}
\]
The multiplicative gate \(M(x)\), for \(x \in \Zp^\ast\), is encoded
as
\[
  M(x) = S^x \cdot H \cdot S^{x^{-1}} \cdot H \cdot S^x \cdot H.
\]
This definition appears both in the README example and in the main
symplectic modules.

\section{Comprehensive Relations}

The relation family in \file{N/Symplectic.agda} is called
\texttt{\_QRel,\_===\_}.  It is indexed by the number of wires and is
used as the axiom relation for the presentation closure.  The
implemented constructors include:

\begin{itemize}
\item one-qupit relations \(S^p = \epsilon\), \(H^4 = \epsilon\),
  \((S H)^3 = \epsilon\), and \(H H S = S H H\);
\item multiplicative-gate laws \(M(x)M(y)=M(xy)\) and
  \(M(x)S = S^{x^2}M(x)\);
\item two-qupit scaling laws for \(M(x)\) through \(CZ\) on either
  side;
\item \(CZ^p=\epsilon\) and commutation of \(CZ\) with lifted
  \(S\)-gates;
\item Selinger-style two- and three-qupit relations named
  \texttt{selinger-c10} through \texttt{selinger-c15};
\item commutation rules saying sufficiently shifted generators commute
  with \(H\), \(S\), and \(CZ\);
\item a congruence-lifting axiom \texttt{cong-up}, written
  \texttt{cong} followed by an up-arrow in the code, which lifts an
  equality on \(n\) wires to one on \(n+1\) wires.
\end{itemize}

The module also proves basic lifting lemmas.  For example, if
\(w \approx v\) under the \(n\)-wire presentation, then
\(w^\uparrow \approx v^\uparrow\) under the \((n+1)\)-wire
presentation.  In the code this is \texttt{lemma-cong-up}, written
with an up-arrow in the actual identifier, proved by induction on the
derivation of \(w \approx v\).

\section{Simplified and Derived Presentations}

The repository distinguishes at least three related presentations.
The comprehensive presentation in \file{N/Symplectic.agda} is
convenient for proofs because it includes relations such as
\(M(x)M(y)=M(xy)\) for arbitrary nonzero \(x,y\).  The simplified
presentation in \file{N/Symplectic-Simplified.agda} is closer to the
paper-facing axiom set.  It is parameterized by a primitive root
\(g \in \Zp^\ast\) and a proof that every nonzero element is a power
of \(g\).  Its one-qupit multiplicative axioms use \(M(g)\) and
powers of \(M(g)\) rather than arbitrary \(M(x)\).

The simplified relation family contains constructors such as:
\[
\begin{aligned}
S^p &= \epsilon,\\
H^2 &= M(-1),\\
M(g)^k &= M(g^k),\\
M(g)S &= S^{g^2}M(g).
\end{aligned}
\]
It also contains the same two- and three-qupit structural relations
needed to reason about \(CZ\), lifted gates, and shifted wires.  This
matches the README's single-ququint example, where the primitive root
is \(2 \in \mathbb{Z}_5^\ast\), and the simplified axiom set derives
the box relation for \(\EBox\).

The derived presentation in \file{N/Symplectic-Derived.agda} exposes
parameterized generators, for example \(H^k\), \(S^k\), and \(CZ^k\),
as generators rather than only as powers of primitive words.  The
module \file{N/Iso-Sym-Derived.agda} defines maps between the basic
and derived presentations:
\[
  f : \Gen_{\mathrm{basic}}(n) \to \Gen_{\mathrm{derived}}(n),
  \qquad
  g : \Gen_{\mathrm{derived}}(n) \to \Word(\Gen_{\mathrm{basic}}(n)).
\]
It then proves well-definedness lemmas for the maps and establishes
isomorphism infrastructure.  The module \file{N/Action-Sym.agda}
uses the composition of these maps to define the Pauli action of the
simplified presentation through the derived one.

\section{Pauli Action}

The semantic side of the development is represented by Pauli vectors.
In \file{N/Pauli.agda}, a one-qupit Pauli is a pair of field elements:
\[
  \mathsf{Pauli1} = \Zp \times \Zp.
\]
An \(n\)-qupit Pauli is a vector of \(n\) such pairs.  The symplectic
form is:
\[
  \mathsf{sform1}((a,b),(c,d)) = -a d + c b,
\]
extended componentwise over vectors.

The module \file{N/Action.agda} defines the action of derived
generators on Pauli vectors.  The action of \(S^k\) maps
\((a,b)\) to \((a,b+a k)\); the action of \(H\) rotates
\((a,b)\) to \((-b,a)\) for exponent \(1\); and \(CZ^k\) updates the
two adjacent components by adding cross terms:
\[
  (a,b),(a',b') \mapsto (a,b+a'k),(a',b'+ak).
\]
Lifted generators act on the tail of the vector, leaving the first
component unchanged.

This Pauli action is used by the normal-form modules.  It lets the
development state and prove facts of the form ``there exists a word in
normal form whose action sends one Pauli vector to another with the
same symplectic invariant.''  The one-qupit normal-form module
\file{N/NF1.agda} is the clearest complete example: it contains
theorems such as \texttt{Theorem-NF1}, \texttt{Theorem-MC}, and
variants specialized to \(pZ\)-like inputs.

\paragraph{Symplectic invariance.}
The action just described is the matrix action \(\rho\) of
Section~\ref{sec:math}: each generator acts by a matrix in
\(\Sp(2n,\Zp)\), and a word acts by the product of the matrices of its
letters, so the whole congruence class of a word has a single symplectic
image.  Every generator preserves \texttt{sform1}, hence so does every
derivable equality; the quantity \(\langle\,\cdot\,,\cdot\,\rangle\) is
the invariant carried along a derivation.  The normal-form theorems are
thus statements about orbits of \(\Sp(2n,\Zp)\) acting on Pauli vectors,
which is also the basis of the efficient classical description of
stabilizer dynamics~\cite{Aaronson2004,deBeaudrap2013}.

\section{Boxes and Normal Forms}

The box syntax is defined in \file{N/Boxes.agda}.  The definitions are
small but heavily indexed:
\[
\begin{aligned}
\ABox &= \{(a,b) \in \Zp^2 \mid (a,b) \ne (0,0)\},\\
\BBox &= \Zp \times \Zp,\\
\DBox &= \Zp \times \Zp,\\
\EBox &= \Zp.
\end{aligned}
\]
The wider box families are indexed by arity.  For example,
\(\LBox(n)\), \(\MBox(n)\), \(\LM(n)\), and \(\NF(n)\) are recursive
families describing normal-form-shaped circuits over \(n\) wires.  A
normal form for \(n+1\) wires consists of a normal form on \(n\) wires
and an \(\LM\) layer on the last wire block:
\[
  \NF(0) = \top,\qquad
  \NF(n+1) = \NF(n) \times \LM(n+1).
\]

The box syntax itself is not a circuit.  The interpretation appears
in \file{N/LM-Sym.agda}, which maps boxes to words over the
symplectic generator set.  Selected definitions are:
\[
\begin{aligned}
[b]^\EBox &= S^{-b},\\
[(0,b)]^\BBox &= Ex \cdot CX'^{\,b},\\
[(0,b)]^\DBox &= Ex \cdot CZ^{-b}.
\end{aligned}
\]
For nonzero first components, the interpretations introduce inverses
in \(\Zp^\ast\).  For instance, a nonzero \(\ABox(a,b)\) is
interpreted using \(a^{-1}\) and \(-b/a\).  The code enforces the
side condition by pattern matching on whether the first component is
zero.

The vector box interpretations are also order-sensitive.  The comment
in \file{N/LM-Sym.agda} states that a vector of \(\BBox\) boxes is in
circuit order, while the corresponding matrix multiplication order is
reversed; the word interpretation accounts for this by recursively
lifting the tail and then composing the current box.

\section{Box Relations}

The public index of box relations is \file{BoxRelations.agda}.  It is
parameterized by \(p-2\) and a proof that \(2+p-2\) is prime.  It
imports the symplectic presentation, the box interpretation, and the
specific proof modules under \file{N/BR}.  The file is organized into
one-, two-, and three-qupit modules.

The one-qupit statements are:
\[
\begin{aligned}
[x]^\ABox \cdot H &\approx dir \cdot [x']^\ABox,\\
[x]^\ABox \cdot S &\approx dir \cdot [x']^\ABox,\\
[b]^\EBox \cdot S &\approx [b-1]^\EBox.
\end{aligned}
\]
The actual direction word and updated box parameter are computed by
functions in \file{N/BR/One/A.agda}; the \(\EBox\) rule is proved in
\file{N/BR/One/E.agda}.

The two-qupit statements are:
\[
\begin{aligned}
[l]^\LBox \cdot CZ &\approx dir \cdot [l']^\LBox,\\
[b]^\BBox \cdot [g] &\approx dir \cdot [b']^\BBox,\\
[d]^\DBox \cdot [g] &\approx S^e \cdot dir^\uparrow \cdot [d']^\DBox.
\end{aligned}
\]
Here the \(\BBox\) rule ranges over the gates \(H^\uparrow\),
\(S^\uparrow\), and \(S\), excluding \(H\) and \(CZ\) by explicit
inequality hypotheses.  The \(\DBox\) rule ranges over \(H\),
\(S^\uparrow\), \(S\), and \(CZ\), excluding \(H^\uparrow\).

The three-qupit statements are:
\[
\begin{aligned}
[vb]^{\vec{\BBox}} \cdot CZ^\uparrow
  &\approx dir^\downarrow \cdot [vb']^{\vec{\BBox}},\\
[b]^\BBox{}^\uparrow \cdot CZ
  &\approx dir \cdot [b]^\BBox{}^\uparrow,\\
[vd]^{\vec{\DBox}} \cdot CZ
  &\approx dir^\uparrow \cdot [vd']^{\vec{\DBox}}.
\end{aligned}
\]
These are delegated to \file{N/BR/Three/BB-CZ.agda},
\file{N/BR/Three/B-CZ.agda}, and
\file{N/BR/Three/DD-CZ.agda}, respectively.

\section{Example: The \(\EBox\) Rule}

The README includes one compact derivation, and it is useful because
it shows the style of the whole project.  The \(\EBox\) interpretation
is \( [b]^\EBox = S^{-b}\).  The box relation says:
\[
  [b]^\EBox \cdot S \approx [b-1]^\EBox.
\]
The proof in \file{N/BR/One/E.agda}, also shown in
\file{doc/Readme.lagda.tex}, proceeds by unfolding the interpretation
and using a lemma that combines powers of \(S\):
\[
  S^{-b} \cdot S \approx S^{-b+1}.
\]
Then finite-field arithmetic rewrites \(-b+1\) into \(-(b-1)\), which
is exactly the exponent used by \([b-1]^\EBox\).

This example illustrates the division of labor in the repository.  A
box theorem is not a custom decision procedure.  It is a typed
calculation in the presentation closure, with proof steps supplied by
previously proved lemmas about powers, modular arithmetic, and word
associativity.

\section{Generic Presentation Machinery}

Although the top-level goal is about Clifford circuits, a significant
part of the repository is generic.  The modules under
\file{Presentation/Construct} define ways to build larger
presentations from smaller ones, including direct products,
semidirect products, and amalgamations.  The modules
\file{Presentation/CosetNF.agda} and
\file{Presentation/Reidemeister-Schreier.agda} formalize coset and
rewriting infrastructure.  The older \file{Presentation/Groups}
directory contains example and exploratory group presentations such
as cyclic groups, symmetric groups, and Clifford one- and two-qupit
fragments.

The Clifford development uses this infrastructure in two ways.  First,
the basic closure and setoid machinery are used directly throughout
the proof scripts.  Second, the normal-form and isomorphism modules
reuse the general idea that a presentation can be related to another
presentation by a well-defined star extension of a generator map.
For a map \(f : X \to \Word Y\), the extension \(f^\ast :
\Word X \to \Word Y\) is implemented by mapping generators to words
and then concatenating.  Well-definedness means that each axiom of
the source presentation maps to a derivable equality in the target
presentation.

\section{Artifact Boundary}

The repository contains unfinished code, visible as \Agda holes
\texttt{\{!!\}} or \texttt{?} in several older or exploratory files.
This is important for an artifact paper: the proof claims should be
scoped to the modules that are intended to constitute the development.
The README itself tells the reader to start from
\file{doc/Readme.pdf} and \file{BoxRelations.agda}; the latter is the
best public index for the box-relation statements.

The distinction is not merely administrative.  Some files such as
\file{N/Completeness.agda}, \file{N/Symplectic-Alternative.agda}, and
older group-presentation experiments contain partially developed
normal-form or alternative-presentation work.  Those modules are
valuable context, but they should not be cited as completed theorems
unless they are closed and type checked in the release artifact.  By
contrast, modules marked \texttt{--safe} and free of holes can be
treated as checked under \Agda's kernel, subject to the standard
library and the explicit postulates or options they import.

For a POPL artifact, the natural packaging step is to provide a small
driver module that imports exactly the intended theorem surface:
\file{BoxRelations.agda}, the simplified-presentation isomorphism
modules, and the finite-field modules used by them.  Such a driver
would make the checked boundary mechanically visible.

\section{Trusted Computing Base}

The trusted base of the development is conventional for an \Agda
formalization:

\begin{itemize}
\item the \Agda type checker;
\item the Agda standard library, tested in the README with Agda 2.8
  and standard library 2.3, and also with Agda 2.7 and standard
  library 2.2;
\item the definitions of the Clifford generators, relation families,
  and box interpretations;
\item the mathematical adequacy claim connecting the formal
  generators and relations to the intended circuit diagrams.
\end{itemize}

The intended derivations themselves do not rely on an external
computer algebra system.  Modular arithmetic is proved inside the
repository using \texttt{Fin}, natural-number modular arithmetic, and
standard-library algebraic structures.  Associativity normalization
is also proved inside the presentation layer, by flattening words to
lists and rebuilding them.

\section{How to Read the Code}

A reader who wants to audit the development can follow this path.

\begin{enumerate}
\item Read \file{Word/Base.agda} and \file{Presentation/Base.agda}.
  These define the syntax and equality relation used everywhere else.
\item Read \file{Zp/ModularArithmetic.agda} enough to understand
  \(\Zp\), nonzero elements, inverse notation, and ring lemmas.
\item Read the first half of \file{N/Symplectic.agda}: generators,
  derived words, \(M(x)\), and the comprehensive relation family.
\item Read \file{N/Boxes.agda} and the beginning of
  \file{N/LM-Sym.agda}: these explain what boxes are and how they
  are interpreted as circuit words.
\item Read \file{BoxRelations.agda}.  It is short and gives the
  theorem statements for the box relations.
\item Follow one proof pointer at a time into \file{N/BR}.  The
  one-qupit \(\EBox\) proof is the easiest starting point; the
  two- and three-qupit proofs rely on more associativity and
  commutation lemmas.
\item Read \file{N/Symplectic-Simplified.agda},
  \file{N/Iso-Sym-Derived.agda}, and \file{N/Action-Sym.agda} for
  the connection between the paper-facing axiom set, derived
  generators, and semantic Pauli actions.
\end{enumerate}

\section{Related Work}

The mathematical context is the theory of stabilizer circuits and the
Clifford group in odd prime dimension.  The symplectic description of
qudit Cliffords used above goes back to the stabilizer
formalism~\cite{Gottesman1997,Gottesman1998} and its higher-dimensional
extensions~\cite{Gottesman1999,Hostens2005,Appleby2005,deBeaudrap2013};
the efficient classical simulation of stabilizer
circuits~\cite{Aaronson2004} and discrete phase-space methods for odd
dimensions~\cite{Gross2006} rely on the same linear-symplectic action.
The splitting of the Clifford-to-symplectic sequence over a finite field
is the Weil representation~\cite{Neuhauser2002}.  Magic-state and
resource-theoretic results single out odd prime dimensions as
especially well behaved~\cite{Veitch2014,Campbell2012,Howard2014}, which
is the regime the present development targets.

Finite presentations of the Clifford group are the direct ancestor of
the relation sets in \file{N/Symplectic.agda}.  Selinger gave generators
and relations for the \(n\)-qubit Clifford operators~\cite{Selinger2015};
the constructors \texttt{selinger-c10}--\texttt{c15} are named after that
style of two- and three-qupit relation.  Farinholt gave an ideal
characterization and presentation of the qudit Clifford
operators~\cite{Farinholt2014}.  Diagrammatic calculi offer an
alternative complete equational theory: the ZX-calculus is complete for
qubit stabilizer quantum mechanics~\cite{CoeckeDuncan2011,Backens2014}
and for the qutrit stabilizer fragment~\cite{Wang2018}.

The proof-engineering context is mechanized reasoning about quantum
programs and about algebraic presentations.  Formal verification of
quantum circuits has been carried out in Coq with
\textsc{qwire}~\cite{Paykin2017} and the \textsc{voqc} verified
optimizer~\cite{Hietala2021}, and via path-sum and related semantic
methods~\cite{Amy2018}.  The present work instead uses
\Agda~\cite{Norell2007,BoveDybjerNorell2009} and its standard
library~\cite{AgdaStdlib}, in the setoid-based style where raw syntax is
kept separate from the congruence that quotients it.  The generic
presentation layer---congruence closure, coset normal forms, and
Reidemeister--Schreier rewriting---follows classical combinatorial and
computational group
theory~\cite{MagnusKarrassSolitar,LyndonSchupp,Johnson1997,Sims1994,HoltEickOBrien}.
The finite-field layer is standard~\cite{LidlNiederreiter,IrelandRosen}.

\section{Conclusion}

\Qupit is a mechanized proof development for Clifford-circuit box
relations over odd prime dimensions.  Its central move is simple and
effective: define circuit syntax as binary words, define equality as
the congruence closure of a presentation, and prove each diagrammatic
rewrite as an \Agda term in that equality.  Around this core, the
repository builds finite-field arithmetic, Clifford generators,
simplified and derived relation sets, Pauli actions, box
interpretations, and a collection of concrete box-relation proof
scripts.

The current code base also records its own research history.  Some
older modules are unfinished, and a POPL artifact should make the
intended checked surface explicit.  The file \file{BoxRelations.agda}
already provides the right shape for that surface: it states the
one-, two-, and three-qupit box relations and points to the modules
where their derivations live.  This paper expands the existing README
into a structured explanation of how those pieces fit together and
what must be checked to support the artifact claims.

\begin{acks}
Omitted for anonymous review.
\end{acks}

\begin{thebibliography}{99}

\bibitem{Agda}
The Agda Team.
\newblock Agda documentation.
\newblock \url{https://agda.readthedocs.io/}.
\newblock Accessed 2026-01-15 in the repository documentation.

\bibitem{AgdaStdlib}
The Agda Standard Library Team.
\newblock Agda standard library documentation, version 2.3.
\newblock \url{https://agda.github.io/agda-stdlib/v2.3/}.
\newblock Accessed 2026-01-15 in the repository documentation.

\bibitem{Gottesman1997}
D.~Gottesman.
\newblock Stabilizer Codes and Quantum Error Correction.
\newblock PhD thesis, California Institute of Technology, 1997.
\newblock arXiv:quant-ph/9705052.

\bibitem{Gottesman1998}
D.~Gottesman.
\newblock The Heisenberg representation of quantum computers.
\newblock In \emph{Group22: Proc. XXII Int. Colloquium on Group Theoretical
  Methods in Physics}, pp.~32--43, 1999.
\newblock arXiv:quant-ph/9807006.

\bibitem{Gottesman1999}
D.~Gottesman.
\newblock Fault-tolerant quantum computation with higher-dimensional systems.
\newblock \emph{Chaos, Solitons \& Fractals}, 10(10):1749--1758, 1999.
\newblock arXiv:quant-ph/9802007.

\bibitem{Aaronson2004}
S.~Aaronson and D.~Gottesman.
\newblock Improved simulation of stabilizer circuits.
\newblock \emph{Physical Review A}, 70:052328, 2004.
\newblock arXiv:quant-ph/0406196.

\bibitem{Hostens2005}
E.~Hostens, J.~Dehaene, and B.~De~Moor.
\newblock Stabilizer states and Clifford operations for systems of arbitrary
  dimensions and modular arithmetic.
\newblock \emph{Physical Review A}, 71:042315, 2005.
\newblock arXiv:quant-ph/0408190.

\bibitem{deBeaudrap2013}
N.~de~Beaudrap.
\newblock A linearized stabilizer formalism for systems of finite dimension.
\newblock \emph{Quantum Information \& Computation}, 13(1--2):73--115, 2013.
\newblock arXiv:1102.3354.

\bibitem{Farinholt2014}
J.~M. Farinholt.
\newblock An ideal characterization of the Clifford operators.
\newblock \emph{Journal of Physics A: Mathematical and Theoretical},
  47(30):305303, 2014.
\newblock arXiv:1307.5087.

\bibitem{Appleby2005}
D.~M. Appleby.
\newblock Symmetric informationally complete-positive operator valued measures
  and the extended Clifford group.
\newblock \emph{Journal of Mathematical Physics}, 46:052107, 2005.
\newblock arXiv:quant-ph/0412001.

\bibitem{Gross2006}
D.~Gross.
\newblock Hudson's theorem for finite-dimensional quantum systems.
\newblock \emph{Journal of Mathematical Physics}, 47:122107, 2006.
\newblock arXiv:quant-ph/0602001.

\bibitem{Neuhauser2002}
M.~Neuhauser.
\newblock An explicit construction of the metaplectic representation over a
  finite field.
\newblock \emph{Journal of Lie Theory}, 12(1):15--30, 2002.

\bibitem{Veitch2014}
V.~Veitch, S.~A.~H. Mousavian, D.~Gottesman, and J.~Emerson.
\newblock The resource theory of stabilizer quantum computation.
\newblock \emph{New Journal of Physics}, 16:013009, 2014.
\newblock arXiv:1307.7171.

\bibitem{Campbell2012}
E.~T. Campbell, H.~Anwar, and D.~E. Browne.
\newblock Magic-state distillation in all prime dimensions using quantum
  Reed--Muller codes.
\newblock \emph{Physical Review X}, 2:041021, 2012.
\newblock arXiv:1205.3104.

\bibitem{Howard2014}
M.~Howard, J.~Wallman, V.~Veitch, and J.~Emerson.
\newblock Contextuality supplies the `magic' for quantum computation.
\newblock \emph{Nature}, 510:351--355, 2014.
\newblock arXiv:1401.4174.

\bibitem{Selinger2015}
P.~Selinger.
\newblock Generators and relations for \(n\)-qubit Clifford operators.
\newblock \emph{Logical Methods in Computer Science}, 11(2):10, 2015.
\newblock arXiv:1310.6813.

\bibitem{CoeckeDuncan2011}
B.~Coecke and R.~Duncan.
\newblock Interacting quantum observables: categorical algebra and
  diagrammatics.
\newblock \emph{New Journal of Physics}, 13:043016, 2011.
\newblock arXiv:0906.4725.

\bibitem{Backens2014}
M.~Backens.
\newblock The ZX-calculus is complete for stabilizer quantum mechanics.
\newblock \emph{New Journal of Physics}, 16:093021, 2014.
\newblock arXiv:1307.7025.

\bibitem{Wang2018}
Q.~Wang.
\newblock Qutrit ZX-calculus is complete for stabilizer quantum mechanics.
\newblock In \emph{Proc. QPL 2018}, EPTCS 287:313--344, 2019.
\newblock arXiv:1803.00696.

\bibitem{Paykin2017}
J.~Paykin, R.~Rand, and S.~Zdancewic.
\newblock \textsc{qwire}: a core language for quantum circuits.
\newblock In \emph{Proc. 44th ACM SIGPLAN Symp. on Principles of Programming
  Languages (POPL)}, pp.~846--858. ACM, 2017.

\bibitem{Hietala2021}
K.~Hietala, R.~Rand, S.-H. Hung, X.~Wu, and M.~Hicks.
\newblock A verified optimizer for quantum circuits.
\newblock \emph{Proceedings of the ACM on Programming Languages},
  5(POPL):37, 2021.
\newblock arXiv:1912.02250.

\bibitem{Amy2018}
M.~Amy.
\newblock Towards large-scale functional verification of universal quantum
  circuits.
\newblock In \emph{Proc. QPL 2018}, EPTCS 287:1--21, 2019.
\newblock arXiv:1805.06908.

\bibitem{Norell2007}
U.~Norell.
\newblock Towards a practical programming language based on dependent type
  theory.
\newblock PhD thesis, Chalmers University of Technology, 2007.

\bibitem{BoveDybjerNorell2009}
A.~Bove, P.~Dybjer, and U.~Norell.
\newblock A brief overview of Agda---a functional language with dependent
  types.
\newblock In \emph{Theorem Proving in Higher Order Logics (TPHOLs)}, LNCS 5674,
  pp.~73--78. Springer, 2009.

\bibitem{MagnusKarrassSolitar}
W.~Magnus, A.~Karrass, and D.~Solitar.
\newblock \emph{Combinatorial Group Theory}.
\newblock Dover, 2nd revised edition, 1976.

\bibitem{LyndonSchupp}
R.~C. Lyndon and P.~E. Schupp.
\newblock \emph{Combinatorial Group Theory}.
\newblock Springer, 1977.

\bibitem{Johnson1997}
D.~L. Johnson.
\newblock \emph{Presentations of Groups}.
\newblock London Mathematical Society Student Texts 15. Cambridge University
  Press, 2nd edition, 1997.

\bibitem{Sims1994}
C.~C. Sims.
\newblock \emph{Computation with Finitely Presented Groups}.
\newblock Encyclopedia of Mathematics and its Applications 48. Cambridge
  University Press, 1994.

\bibitem{HoltEickOBrien}
D.~F. Holt, B.~Eick, and E.~A. O'Brien.
\newblock \emph{Handbook of Computational Group Theory}.
\newblock Chapman \& Hall/CRC, 2005.

\bibitem{LidlNiederreiter}
R.~Lidl and H.~Niederreiter.
\newblock \emph{Finite Fields}.
\newblock Encyclopedia of Mathematics and its Applications 20. Cambridge
  University Press, 2nd edition, 1997.

\bibitem{IrelandRosen}
K.~Ireland and M.~Rosen.
\newblock \emph{A Classical Introduction to Modern Number Theory}.
\newblock Graduate Texts in Mathematics 84. Springer, 2nd edition, 1990.

\end{thebibliography}

\end{document}
